\chapter{导出、保存与恢复}\label{ch:export}
\section{选择正确的出口}
\par\noindent\begin{tabularx}{\linewidth}{@{}P{33mm}P{39mm}Y@{}}
\toprule
目的 & 入口或文件 & 需要知道的区别\\\midrule
继续编辑 / 跨机传递 & 保存项目，\code{.inkdoc} & 包含内容、设置和内嵌资源；优先作为工作母版。\\
只交换文字 & 保存 Markdown / 文本 & 不带完整字形、评价和样式。\\
查看整页外观 & 更多导出选项 → 导出参考页面… & 输出 SVG/PNG；位图参考页不能当作机器文件。\\
发送写字机 & 顶部导出 SVG… & 完整文档逐页导出，仅可书写路径。\\
分析字形质量 & 更多导出选项 → 导出评价 CSV… & 输出评价和来源，不修改生成器。\\
\bottomrule
\end{tabularx}

页面与机器导出选择尚不存在的新目录，避免覆盖已有结果。多页文件应按页序核对；不要只传第一页。写字机输出以 \code{page-001.svg} 等命名，\code{plotter.json} 记录路径来源、优化距离及相关说明。

\section{写字机 SVG 的含义}
机器页面中的绘制元素仅有 \code{path}，路径指令仅有绝对 \code{M} 和 \code{L}。字体、文字元素、图片和变换矩阵不会留给设备解释。\code{M} 定位到笔画起点，\code{L} 沿折线书写；具体抬笔、速度和力度由设备软件决定。

导出前依次确认：最新预览已经完成；没有阻断错误；字符与图形都具有可书写路径；页面与模板模式正确。缺字、越界和不可支持内容不能通过“忽略提醒”变成完整成品。

含普通位图时，点击导出会说明这些图片不会进入 SVG，包括尚未载入的本地 PNG/JPEG。弹窗分别显示两类数量，展开详情可查看文件名。确认后保留当前图片占位与图题，导出其余文字及受支持的单线图路径；取消不产生输出。原文档及资源不变，\code{plotter.json} 记录原因、页码、源码身份和占位；缺失图片哈希为 null。预览和适配报告继续显示资源缺失错误，整份文档仍未完整可导出。缺失 SVG、未知类型、非法资源、缺字、公式错误与越框仍须修正。需要图片本身参与书写时，应准备经过确认的单线图资源。

\section{比例校正与输出画布}
先严格校验未缩放排版，再以左上角为原点同步缩放输出坐标、SVG 物理宽高及 viewBox，笔宽保持毫米值。例如横纵均为 132.293\% 时，输出画布宽高都乘以 1.32293。文档纸张、预览、分页及冻结选样不改变。选定页独立变换，文件沿用原页码；\code{plotter.json} 记录原始及输出尺寸与变换策略。

缩放前的越界等阻断错误不能通过缩小或筛选规避。缩放后仍校验画布边界、路径及单边 2000 mm 上限；笔宽不随缩小改变，因此贴边笔迹缩小后仍可能被拒绝。定位短线坐标随比例变换，笔宽仍为 0.01 mm。

\section{优化空行程}
\entry{导出 SVG… → 优化路径（允许反转笔画执行方向）}
默认开启。开启后，程序调整路径执行顺序，并允许反向执行某条路径，以减少抬笔移动距离；反转不是另一个可取消的选项。排序和反转不改变路径形状及纸面位置，也不改写原始采样文件。

模板规则线、固定字段和正文等执行区间仍按既定顺序处理，不会为了缩短空行程任意混在一起。清单中的距离表示几何估计，不能直接当作实际节省的时间；落笔方向对具体笔具的影响仍以试写为准。

第一次使用某组样本或笔具时，可以分别导出开启/关闭优化的两组文件比较。两组输出用不同目录，并保留当时的项目版本。

\section{保存项目与多人协作}
\ui{保存项目}保留源码、样式、资源、模板、绑定、选样和评价。\ui{另存项目…}用于试验布局或保留不同版本；\ui{保存 Markdown / 文本…}只输出源码。

从 0.1.82 起，Ctrl+S 在后台保存 Markdown 项目。保存期间可以继续编辑。无改动时跳过写盘；若已有新改动，多次保存请求会合并，当前保存完成后再保存最新内容。

同一项目在另一窗口或另一台机器被修改后，保存可能报告冲突。先另存自己的副本，再比较内容，不反复强行覆盖。编辑器的冲突检测不是多人实时协作或合并系统。

正式更新已有项目时会保留 \code{.previous.inkdoc} 副本。跨机传递后，应在目标机打开项目并重新查看预览；不要只看文件存在就认为资源齐全。复制原始 Markdown 时，还需另外提供字形与图像资源。

\section{打开恢复副本}
\entry{打开 → 打开恢复副本…}
独立应用按会话写自动恢复副本；macOS 默认位置为：
\par\noindent\begin{minipage}{\linewidth}
\begin{Source}
~/Library/Application Support/InkStruct/editor_recovery
\end{Source}
\end{minipage}\par\medskip
源码运行方式的恢复目录是启动工作目录下的 \code{artifacts/editor_recovery/}。每个会话独立保存，重开同一项目或多个窗口不会共用一个草稿文件。

恢复后项目处于未保存状态，核对正文与设置，再选择正式位置保存；恢复副本不会自动覆盖原项目。遇到损坏项目，打开失败应保留当前编辑内容，可继续另存。超限或缺字文档可以保存草稿，但仍需修正后才能导出成品。

\section{交付一份可继续处理的文档}
建议交付原生项目、机器输出目录，以及简短的纸张/笔具/模板模式说明。模板需要复用时另附模板 JSON。若只交付 SVG，接收者可书写，但难以重新调整正文、字体样本或模板绑定。

\section{后台 SVG 导出（0.1.38）}
确认排版在输出位置选择前冻结，导出弹窗中的页码、表格内容、路径优化、横纵比例和定位点设置共同作用于同一份冻结快照。后台期间可继续编辑和滚动，不会改变本次输出；每个窗口同时只运行一次导出。

多页路径优化最多使用四个进程，并受处理器和可用内存限制；单页、小任务及缓存命中沿用串行或缓存路径。进程失败会对未完成页串行重算，不放宽几何和来源检查。

进度窗口可取消，关闭进度窗口也请求取消。取消后清理临时文件，不发布半成品；最终发布已完成时保留成功输出。退出应用须先等待导出取消和清理完成，再次关闭。已有输出目录不能覆盖。
